\documentclass[11pt,a4paper]{article}
\usepackage[utf8]{inputenc}
\usepackage[margin=2.5cm]{geometry}
\usepackage{amsmath,amssymb}
\usepackage{booktabs}
\usepackage{xcolor}
\usepackage{colortbl}
\usepackage{hyperref}

\definecolor{primary}{RGB}{15, 76, 129}
\definecolor{secondary}{RGB}{240, 244, 248}
\definecolor{accent}{RGB}{220, 53, 69}

\hypersetup{
    colorlinks=true,
    linkcolor=primary,
    urlcolor=primary
}

\title{\textbf{\Large Lógica e Sistemas Digitais --- TLAB1}\\
\large Controlo de Acessos: Tabelas de Verdade, Mapas de Karnaugh e VHDL}
\author{\textbf{Instituto Superior de Engenharia de Lisboa (ISEL)}}
\date{Ano Letivo 2026/2027}

\begin{document}

\maketitle

\section{Introdução}
Este documento apresenta o desenvolvimento do \textbf{Trabalho Laboratorial 1 (TLAB1)} da unidade curricular de \textbf{Lógica e Sistemas Digitais}. O objetivo é o projeto, simplificação e implementação em VHDL de um sistema combinatório de Controlo de Acessos para 16 funcionários e 4 departamentos.

\section{Tabela de Verdade Global}
A Tabela~\ref{tab:verdade} mapeia o código binário do funcionário ($A_3 A_2 A_1 A_0$), o respetivo mintermo ($m_0$ a $m_{15}$) e a autorização de acesso para cada um dos departamentos ($D_3, D_2, D_1, D_0$).

\begin{table}[h!]
\centering
\caption{Tabela de Verdade Completa do Sistema.}
\label{tab:verdade}
\vspace{0.2cm}
\begin{tabular}{ccccccccccc}
\toprule
\textbf{Func.} & \textbf{$A_3$} & \textbf{$A_2$} & \textbf{$A_1$} & \textbf{$A_0$} & \textbf{Mintermo} & \textbf{Departamentos} & \textbf{$D_3$} & \textbf{$D_2$} & \textbf{$D_1$} & \textbf{$D_0$} \\
\midrule
\rowcolor{secondary} 0 & 0 & 0 & 0 & 0 & $m_0$ & $D_0$ & 0 & 0 & 0 & 1 \\
1 & 0 & 0 & 0 & 1 & $m_1$ & $D_2$ & 0 & 1 & 0 & 0 \\
\rowcolor{secondary} 2 & 0 & 0 & 1 & 0 & $m_2$ & $D_1$ & 0 & 0 & 1 & 0 \\
3 & 0 & 0 & 1 & 1 & $m_3$ & $D_3$ & 1 & 0 & 0 & 0 \\
\rowcolor{secondary} 4 & 0 & 1 & 0 & 0 & $m_4$ & $D_1, D_2$ & 0 & 1 & 1 & 0 \\
5 & 0 & 1 & 0 & 1 & $m_5$ & $D_2, D_0$ & 0 & 1 & 0 & 1 \\
\rowcolor{secondary} 6 & 0 & 1 & 1 & 0 & $m_6$ & $D_0$ & 0 & 0 & 0 & 1 \\
7 & 0 & 1 & 1 & 1 & $m_7$ & $D_2$ & 0 & 1 & 0 & 0 \\
\rowcolor{secondary} 8 & 1 & 0 & 0 & 0 & $m_8$ & $D_3$ & 1 & 0 & 0 & 0 \\
9 & 1 & 0 & 0 & 1 & $m_9$ & $D_3$ & 1 & 0 & 0 & 0 \\
\rowcolor{secondary} 10 & 1 & 0 & 1 & 0 & $m_{10}$ & $D_0$ & 0 & 0 & 0 & 1 \\
11 & 1 & 0 & 1 & 1 & $m_{11}$ & $D_2, D_0$ & 0 & 1 & 0 & 1 \\
\rowcolor{secondary} 12 & 1 & 1 & 0 & 0 & $m_{12}$ & $D_2, D_0$ & 0 & 1 & 0 & 1 \\
13 & 1 & 1 & 0 & 1 & $m_{13}$ & $D_1$ & 0 & 0 & 1 & 0 \\
\rowcolor{secondary} 14 & 1 & 1 & 1 & 0 & $m_{14}$ & $D_2, D_1$ & 0 & 1 & 1 & 0 \\
15 & 1 & 1 & 1 & 1 & $m_{15}$ & $D_1, D_3$ & 1 & 0 & 1 & 0 \\
\bottomrule
\end{tabular}
\end{table}

\newpage

\section{\texorpdfstring{Mapas de Karnaugh ($4 \times 4$)}{Mapas de Karnaugh (4x4)}}

A ordenação das linhas ($A_3 A_2$) e colunas ($A_1 A_0$) segue o Código de Gray ($00, 01, 11, 10$).

\subsection{Grelha de Referência dos Mintermos}
\begin{table}[h!]
\centering
\begin{tabular}{|c|c|c|c|c|}
\hline
\rowcolor{primary!15} \textbf{$A_3 A_2 \backslash A_1 A_0$} & \textbf{00} & \textbf{01} & \textbf{11} & \textbf{10} \\ \hline
\textbf{00} & $m_0$ & $m_1$ & $m_3$ & $m_2$ \\ \hline
\textbf{01} & $m_4$ & $m_5$ & $m_7$ & $m_6$ \\ \hline
\textbf{11} & $m_{12}$ & $m_{13}$ & $m_{15}$ & $m_{14}$ \\ \hline
\textbf{10} & $m_8$ & $m_9$ & $m_{11}$ & $m_{10}$ \\ \hline
\end{tabular}
\end{table}

\subsection{\texorpdfstring{Mapa de Karnaugh e Equação --- Departamento $D_0$}{Mapa D0}}
Mintermos ativos: $\{m_0, m_5, m_6, m_{10}, m_{11}, m_{12}\}$

\begin{table}[h!]
\centering
\begin{tabular}{|c|c|c|c|c|}
\hline
\rowcolor{secondary} \textbf{$A_3 A_2 \backslash A_1 A_0$} & \textbf{00} & \textbf{01} & \textbf{11} & \textbf{10} \\ \hline
\textbf{00} & \textbf{\textcolor{accent}{1}} & 0 & 0 & 0 \\ \hline
\textbf{01} & 0 & \textbf{\textcolor{accent}{1}} & 0 & \textbf{\textcolor{accent}{1}} \\ \hline
\textbf{11} & \textbf{\textcolor{accent}{1}} & 0 & 0 & 0 \\ \hline
\textbf{10} & 0 & 0 & \textbf{\textcolor{accent}{1}} & \textbf{\textcolor{accent}{1}} \\ \hline
\end{tabular}
\end{table}

\begin{equation}
D_0 = A_3 \cdot \overline{A_2} \cdot A_1 + A_2 \cdot \overline{A_1} \cdot A_0 + \overline{A_3} \cdot \overline{A_2} \cdot \overline{A_1} \cdot \overline{A_0} + \overline{A_3} \cdot A_2 \cdot A_1 \cdot \overline{A_0} + A_3 \cdot A_2 \cdot \overline{A_1} \cdot \overline{A_0}
\end{equation}

\subsection{\texorpdfstring{Mapa de Karnaugh e Equação --- Departamento $D_1$}{Mapa D1}}
Mintermos ativos: $\{m_2, m_4, m_{13}, m_{14}, m_{15}\}$

\begin{table}[h!]
\centering
\begin{tabular}{|c|c|c|c|c|}
\hline
\rowcolor{secondary} \textbf{$A_3 A_2 \backslash A_1 A_0$} & \textbf{00} & \textbf{01} & \textbf{11} & \textbf{10} \\ \hline
\textbf{00} & 0 & 0 & 0 & \textbf{\textcolor{accent}{1}} \\ \hline
\textbf{01} & \textbf{\textcolor{accent}{1}} & 0 & 0 & 0 \\ \hline
\textbf{11} & 0 & \textbf{\textcolor{accent}{1}} & \textbf{\textcolor{accent}{1}} & \textbf{\textcolor{accent}{1}} \\ \hline
\textbf{10} & 0 & 0 & 0 & 0 \\ \hline
\end{tabular}
\end{table}

\begin{equation}
D_1 = A_3 \cdot A_2 \cdot A_0 + A_3 \cdot A_2 \cdot A_1 + \overline{A_3} \cdot \overline{A_2} \cdot A_1 \cdot \overline{A_0} + \overline{A_3} \cdot A_2 \cdot \overline{A_1} \cdot \overline{A_0}
\end{equation}

\subsection{\texorpdfstring{Mapa de Karnaugh e Equação --- Departamento $D_2$}{Mapa D2}}
Mintermos ativos: $\{m_1, m_4, m_5, m_7, m_{11}, m_{12}, m_{14}\}$

\begin{table}[h!]
\centering
\begin{tabular}{|c|c|c|c|c|}
\hline
\rowcolor{secondary} \textbf{$A_3 A_2 \backslash A_1 A_0$} & \textbf{00} & \textbf{01} & \textbf{11} & \textbf{10} \\ \hline
\textbf{00} & 0 & \textbf{\textcolor{accent}{1}} & 0 & 0 \\ \hline
\textbf{01} & \textbf{\textcolor{accent}{1}} & \textbf{\textcolor{accent}{1}} & \textbf{\textcolor{accent}{1}} & 0 \\ \hline
\textbf{11} & \textbf{\textcolor{accent}{1}} & 0 & 0 & \textbf{\textcolor{accent}{1}} \\ \hline
\textbf{10} & 0 & 0 & \textbf{\textcolor{accent}{1}} & 0 \\ \hline
\end{tabular}
\end{table}

\begin{equation}
D_2 = A_2 \cdot \overline{A_1} \cdot \overline{A_0} + A_3 \cdot A_2 \cdot \overline{A_0} + \overline{A_3} \cdot \overline{A_1} \cdot A_0 + \overline{A_3} \cdot A_2 \cdot A_0 + A_3 \cdot \overline{A_2} \cdot A_1 \cdot A_0
\end{equation}

\subsection{\texorpdfstring{Mapa de Karnaugh e Equação --- Departamento $D_3$}{Mapa D3}}
Mintermos ativos: $\{m_3, m_8, m_9, m_{15}\}$

\begin{table}[h!]
\centering
\begin{tabular}{|c|c|c|c|c|}
\hline
\rowcolor{secondary} \textbf{$A_3 A_2 \backslash A_1 A_0$} & \textbf{00} & \textbf{01} & \textbf{11} & \textbf{10} \\ \hline
\textbf{00} & 0 & 0 & \textbf{\textcolor{accent}{1}} & 0 \\ \hline
\textbf{01} & 0 & 0 & 0 & 0 \\ \hline
\textbf{11} & 0 & 0 & \textbf{\textcolor{accent}{1}} & 0 \\ \hline
\textbf{10} & \textbf{\textcolor{accent}{1}} & \textbf{\textcolor{accent}{1}} & 0 & 0 \\ \hline
\end{tabular}
\end{table}

\begin{equation}
D_3 = A_3 \cdot \overline{A_2} \cdot \overline{A_1} + \overline{A_3} \cdot \overline{A_2} \cdot A_1 \cdot A_0 + A_3 \cdot A_2 \cdot A_1 \cdot A_0
\end{equation}

\section{Bloco Verifica (Sinal de Saída $R$)}
O bloco de verificação compara os departamentos autorizados com o departamento ativo no sistema ($S_1 S_0$):
\begin{equation}
R = D_0 \cdot \overline{S_1} \cdot \overline{S_0} + D_1 \cdot \overline{S_1} \cdot S_0 + D_2 \cdot S_1 \cdot \overline{S_0} + D_3 \cdot S_1 \cdot S_0
\end{equation}

\end{document}
